% TEMPLATE-MILC-v3.tex - plantilla maestra UNICA de las guias MILC v3.
%
% La usan las tres familias de guias, cada una con su driver:
%   · Grados 6-11  -> scripts/build-guias-g11.py
%   · Bebras       -> scripts/build-guias-bebras.py
%   · Semillero    -> scripts/build-guias-semillero.py
%
% Antes habia tres copias casi identicas (~400 lineas iguales, 6 distintas):
% cada mejora habia que aplicarla tres veces, y en la practica no se hacia
% (el motor de recursos vivia solo en dos de ellas). Lo que varia entre
% familias esta parametrizado en 6 placeholders que cada driver llena
% (se nombran aqui SIN su sintaxis real: el builder reemplaza por texto plano,
% tambien dentro de los comentarios, y un valor multilinea rompe el preambulo):
%
%   HEADER_IZQ        encabezado izquierdo de pagina
%   HEADER_DER        encabezado derecho de pagina
%   PORTADA_ETIQUETA  rotulo sobre el numero grande (GRADO/MOMENTO/MODULO)
%   PORTADA_SERIE     linea de serie de la portada
%   PORTADA_ID        identificador de la guia en la portada
%   PORTADA_CIRCULO   el circulito de grado (°), o vacio si no aplica
%
% El resto de claves las documenta content/guias/_SCHEMA.md. El script valida
% que no queden placeholders sin reemplazar antes de escribir el archivo final.

\documentclass[11pt,letterpaper]{article}
\usepackage[margin=2.0cm]{geometry}
\usepackage[table]{xcolor}
\usepackage{fontspec}
\usepackage[spanish]{babel}
\usepackage{tikz}
% Librerías TikZ para los diagramas de las guías (bloque `recursos:`):
%  · babel  → babel en español vuelve activos caracteres como «>», y TikZ los usa
%             en sus opciones (>=stealth, ->). Sin esta librería, cualquier
%             diagrama con flechas revienta con «\language@active@arg> has an
%             extra }». Debe ir ANTES de que se dibuje nada.
%  · positioning        → sintaxis «below left=2pt and 3pt».
%  · decorations.pathreplacing → llaves ({brace}) para señalar rangos.
\usetikzlibrary{babel,positioning,decorations.pathreplacing}
\usepackage{graphicx}
% Circuitos: el trabajo de electrónica se hace hoy con micro:bit y sus sensores
% integrados (MakeCode, sin cableado), así que no se necesitan esquemáticos.
% Si algún día entran montajes con protoboard, basta descomentar esta línea:
% los diagramas .tex/.tikz declarados en `recursos:` ya se incrustan solos.
% \usepackage{circuitikz}
\usepackage[most]{tcolorbox}
\usepackage{tabularx}
\usepackage{array}
\usepackage{fancyhdr}
\usepackage{titlesec}
\usepackage[document]{ragged2e}  % bandera a la derecha en todo el cuerpo
\usepackage{enumitem}
\usepackage{microtype}

% Helvetica Neue no trae la flecha U+2192, asi que cada «→» escrito literalmente
% en el YAML se compilaba a NADA, en silencio (462 flechas en 61 guias: rutas del
% tipo «paleta → tipografia → lockup» salian rotas). Mapeamos el caracter a su
% version matematica, que si tiene glifo. Asi el autor puede seguir escribiendo
% «→» comodamente en el YAML.
\usepackage{newunicodechar}
\newunicodechar{→}{\ensuremath{\rightarrow}}
% Mismo problema con los simbolos de marca y vinetas: el «✓» de «una palomita ✓»
% salia como cuadrito vacio en el PDF de todas las guias que lo usaban.
\usepackage{pifont}
\newunicodechar{✓}{\ding{51}}
\newunicodechar{✗}{\ding{55}}
\newunicodechar{✔}{\ding{52}}
\newunicodechar{•}{\textbullet}
\newunicodechar{≥}{\ensuremath{\geq}}
\newunicodechar{≤}{\ensuremath{\leq}}

\setmainfont{Helvetica Neue}
\setsansfont{Helvetica Neue}
\setlength{\parindent}{0pt}
\setlength{\parskip}{6pt}
% Sin particion de palabras: en 6.º-7.º una palabra cortada por guion al final
% de linea («Comuni-cacion») frena la lectura mucho mas que una linea corta.
\hyphenpenalty=10000
\exhyphenpenalty=10000
\pretolerance=2000
\tolerance=3000
\setlength{\emergencystretch}{3em}
\linespread{1.10}
\renewcommand{\arraystretch}{1.25}
\setlist{leftmargin=5mm,itemsep=1mm,topsep=1mm}
\newcolumntype{Y}{>{\RaggedRight\arraybackslash}X}

\definecolor{milcMagenta}{HTML}{E6007E}
\definecolor{milcVino}{HTML}{5A0038}
\definecolor{milcTurquesa}{HTML}{0093A5}
\definecolor{milcVerde}{HTML}{58A923}
\definecolor{milcMostaza}{HTML}{E5B400}
\definecolor{milcOcre}{HTML}{B86600}
\definecolor{milcNegro}{HTML}{241816}
\definecolor{milcGris}{HTML}{F7F7F4}
\definecolor{milcLinea}{HTML}{E8E2DD}
\definecolor{milcGradePrimary}{HTML}{84CC16}
\definecolor{milcGradeDark}{HTML}{4D7C0F}
\definecolor{milcGradeSoft}{HTML}{ECFCCB}
\definecolor{milcGradeAccent}{HTML}{CFE7FF}
\definecolor{milcGradeText}{HTML}{FFFFFF}
\color{milcNegro}

\pagestyle{fancy}
\fancyhf{}
\lhead{\small Tecnología e Informática · Grado 7}
\rhead{\small Guía 12 · MILC v3}
\cfoot{\small\thepage}
\renewcommand{\headrulewidth}{0pt}

\newtcolorbox{titlebox}[1]{
  enhanced,colback=#1,colframe=#1,arc=7mm,boxrule=0pt,
  left=7mm,right=7mm,top=3mm,bottom=3mm,
  before skip=8pt,after skip=8pt,
  fontupper=\sffamily\bfseries\Large\color{white}
}

\newtcolorbox{softbox}[3]{
  enhanced,colback=#2,colframe=#1,borderline west={4pt}{0pt}{#1},
  arc=2mm,boxrule=.4pt,left=4mm,right=4mm,top=3mm,bottom=3mm,
  before skip=6pt,after skip=8pt,title={#3},coltitle=white,
  colbacktitle=#1,fonttitle=\sffamily\bfseries,fontupper=\RaggedRight,
  attach boxed title to top left={xshift=3mm,yshift=-2mm},
  boxed title style={arc=2mm, boxrule=0pt}
}

\newtcolorbox{drawbox}[2]{
  enhanced,colback=white,colframe=#1,arc=4mm,boxrule=1pt,
  left=5mm,right=5mm,top=4mm,bottom=4mm,before skip=8pt,after skip=8pt,
  title={#2},coltitle=white,colbacktitle=#1,fonttitle=\sffamily\bfseries,
  fontupper=\RaggedRight,
  attach boxed title to top left={xshift=4mm,yshift=-2mm},
  boxed title style={arc=3mm, boxrule=0pt}
}

\newtcolorbox{infoband}[1]{
  enhanced,colback=#1!8,colframe=#1!50,arc=2mm,boxrule=.5pt,
  left=4mm,right=4mm,top=2mm,bottom=2mm,before skip=4pt,after skip=4pt,
  fontupper=\small\RaggedRight
}

% Puente narrativo entre secciones (contrato editorial Conéctate).
% Da continuidad: "Acabas de X. Ahora vas a Y porque Z."
% Visualmente: banda izquierda en color del grado, texto en itálica pequeña.
\newtcolorbox{puentebox}{
  enhanced,colback=milcGradeSoft,colframe=milcGradeDark,
  borderline west={3pt}{0pt}{milcGradeDark},
  arc=0mm,boxrule=0pt,left=5mm,right=4mm,top=2.5mm,bottom=2.5mm,
  before skip=4pt,after skip=8pt,
  fontupper=\itshape\small\color{milcNegro}\RaggedRight
}
% La flecha va en modo matematico a proposito: Helvetica Neue NO tiene el glifo
% U+2192, asi que \textrightarrow se compilaba a NADA («Missing character: There
% is no → in font Helvetica Neue») y el puente salia sin flecha. En modo
% matematico la flecha viene de la fuente matematica y si se dibuja.
\newcommand{\puente}[1]{%
  \begin{puentebox}%
    \textcolor{milcGradeDark}{$\rightarrow$}\hspace{2mm}#1%
  \end{puentebox}%
}

\newcommand{\linead}[1]{\noindent\color{milcLinea}\rule{#1}{0.5pt}\color{milcNegro}}
\newcommand{\respuesta}[1]{\par\vspace{2mm}\foreach \n in {1,...,#1}{\noindent\color{milcLinea}\rule{\linewidth}{0.5pt}\par\vspace{5mm}}\color{milcNegro}}
\newcommand{\checkbox}{\textcolor{milcGradeDark}{\fbox{\phantom{x}}}\hspace{5pt}}

\newcommand{\phasecard}[4]{
\begin{tcolorbox}[enhanced,colback=#3,colframe=#2,arc=3mm,boxrule=.5pt,height=3.05cm,valign=center,halign=center,left=1.5mm,right=1.5mm,top=2mm,bottom=2mm]
{\sffamily\bfseries\fontsize{22}{24}\selectfont\textcolor{#2}{#1}}\par
{\sffamily\bfseries\small #4}
\end{tcolorbox}
}

% ── Piezas del rediseno 2026 (legibilidad 6.º-11.º) ───────────────────
% Banda de estacion: reemplaza el titlebox macizo. Titulo a la izquierda,
% dato practico (tiempo/modalidad) a la derecha, en una sola linea.
\newcommand{\milcestacion}[2]{%
\begin{tcolorbox}[enhanced,colback=#1,colframe=#1,arc=2mm,boxrule=0pt,
  left=5mm,right=5mm,top=2.6mm,bottom=2.6mm,before skip=11pt,after skip=7pt]
{\sffamily\bfseries\fontsize{15}{18}\selectfont\color{white}#2}
\end{tcolorbox}}

% Tarjeta de paso numerada: sustituye las tablas «Pilar / Aplicacion», que
% eran formato de consulta docente, no de estudiante. El numero va en un
% circulo sobre el borde para que la secuencia se lea de un vistazo.
\newcommand{\milcpaso}[3]{%
\begin{tcolorbox}[enhanced,colback=white,colframe=milcLinea,arc=1.5mm,boxrule=.9pt,
  left=14mm,right=4mm,top=3mm,bottom=3mm,before skip=4pt,after skip=4pt,
  fontupper=\RaggedRight,
  overlay={\node[anchor=north west,circle,fill=#2,text=white,inner sep=0pt,
    minimum size=7.4mm,font=\sffamily\bfseries\small]
    at ([xshift=3.6mm,yshift=-3.2mm]frame.north west) {#1};}]
#3
\end{tcolorbox}}

% Casilla marcable de tamano util con lapiz (la anterior era \fbox{\phantom{x}},
% de alto variable segun el texto que la seguia).
\newcommand{\milcmarca}{\framebox[3.8mm]{\rule{0pt}{3.4mm}}}

% Fila marcable de lista suelta (checks de la estacion 1).
\newcommand{\milccheckrow}[1]{%
\par\vspace{1.8mm}\noindent
\makebox[6.4mm][l]{\raisebox{-.55mm}{\textcolor{milcNegro}{\milcmarca}}}%
\parbox[t]{\dimexpr\linewidth-6.4mm\relax}{\RaggedRight #1}\par}

% Celda marcable dentro de la rubrica (tabla de autoevaluacion). Misma
% sangria francesa, pero pensada para vivir dentro de una columna X.
\newcommand{\milcopcion}[1]{%
\makebox[5.2mm][l]{\raisebox{-.55mm}{\textcolor{milcNegro}{\milcmarca}}}%
\parbox[t]{\dimexpr\linewidth-5.2mm\relax}{\RaggedRight #1}}

% ── Recursos visuales de la guía ──────────────────────────────────────
% Declarados en el bloque `recursos:` del YAML y emitidos por el builder.
% \guiaFigura   → imagen rasterizada o PDF (foto, carta celeste, montaje).
% \guiaDiagrama → fuente TikZ/circuitikz (.tex) incrustada como vector:
%                 es la vía para circuitos y esquemáticos editables.
% #1 = ruta relativa al .tex · #2 = pie (puede ir vacío)
\newcommand{\guiaFigura}[2]{%
\begin{center}
\includegraphics[width=0.88\linewidth,height=0.38\textheight,keepaspectratio]{#1}\\[1.5mm]
{\footnotesize\itshape\textcolor{milcNegro!75}{#2}}
\end{center}
}

\newcommand{\guiaDiagrama}[2]{%
\begin{center}
\input{#1}\\[1.5mm]
{\footnotesize\itshape\textcolor{milcNegro!75}{#2}}
\end{center}
}

\begin{document}
\thispagestyle{empty}

% ── Portada ───────────────────────────────────────────────────────────
% Antes: un rectangulo de color a pagina completa. Se veia bien en pantalla,
% pero en la fotocopiadora del colegio se comia el toner y salia gris sucio.
% Ahora la banda ocupa el tercio superior y el resto va en blanco.
\begin{tikzpicture}[remember picture,overlay]
  \path[fill=milcGradePrimary]
    (current page.north west) rectangle ([yshift=-10.6cm]current page.north east);

  \node[anchor=north west,text=milcGradeText] at ([xshift=2.0cm,yshift=-1.55cm]current page.north west)
    {\sffamily\bfseries\fontsize{12}{14}\selectfont GRADO};
  \node[anchor=north west,text=milcGradeText,opacity=.85] at ([xshift=2.0cm,yshift=-2.35cm]current page.north west)
    {\sffamily\bfseries\fontsize{8.5}{10}\selectfont SERIE GUÍAS · TECNOLOGÍA E INFORMÁTICA};
  \node[anchor=north east,text=milcGradeText,opacity=.85,align=right] at ([xshift=-2.0cm,yshift=-1.55cm]current page.north east)
    {\sffamily\bfseries\fontsize{9}{12}\selectfont I.E. Sor María Juliana\\Cartago, Valle del Cauca};

  \node[anchor=west,text=milcGradeText] (gradonum) at ([xshift=1.85cm,yshift=-7.1cm]current page.north west)
    {\sffamily\bfseries\fontsize{112}{112}\selectfont 7};
  % El circulito de grado (°) solo aplica a las guias de grado («11°»); en
  % Bebras y Semillero el numero grande es un momento/modulo, no un grado.
  % Se posiciona relativo al nodo `gradonum`, no en coordenadas absolutas.
  \draw[milcGradeText,line width=1.6pt]
    ([xshift=2.0mm,yshift=-4.5mm]gradonum.north east) circle (.15cm);

  \node[anchor=west,text=milcGradeText,text width=11.4cm,align=left] at ([xshift=1.5cm,yshift=0.5cm]gradonum.east)
    {
     {\sffamily\bfseries\fontsize{9}{11}\selectfont\MakeUppercase{Período 2 · Algoritmia y pensamiento computacional}}\\[3mm]
     {\sffamily\bfseries\fontsize{21}{25}\selectfont\setlength{\baselineskip}{27pt}¿Qué es un\\[4pt]algoritmo? ---\\[4pt]la receta\\[4pt]paso a paso}\\[3.5mm]
     {\sffamily\bfseries\fontsize{15}{18}\selectfont Guía 12-7-TIC}
    };
\end{tikzpicture}

\vspace*{9.4cm}
\vfill

{\sffamily\bfseries\fontsize{8}{10}\selectfont\textcolor{milcGradeDark}{LO QUE TE LLEVAS HOY}}

\begin{tcolorbox}[enhanced,colback=white,colframe=milcNegro,arc=2mm,boxrule=1.6pt,
  left=5mm,right=5mm,top=4mm,bottom=4mm,before skip=3pt,after skip=12pt,
  fontupper=\RaggedRight\fontsize{11}{15.5}\selectfont]
Un \textbf{algoritmo escrito en el cuaderno} de algo cotidiano que tú sepas hacer (preparar arepas, organizar el morral, lavarte los dientes, hacer una llamada, jugar un juego, etc.). El algoritmo tiene: \textbf{(a) nombre claro}; \textbf{(b) entrada} (qué necesitas para empezar); \textbf{(c) pasos numerados} (mínimo 6); \textbf{(d) salida} (qué resultado se obtiene); \textbf{(e) 5 características} verificadas (claro, finito, determinista, eficaz, ordenado). Más \textbf{una autoevaluación} de cuál característica te costó más cumplir.
\end{tcolorbox}

\begin{tcolorbox}[enhanced,colback=milcGris,colframe=milcGris,arc=2mm,boxrule=0pt,
  left=5mm,right=5mm,top=3.5mm,bottom=3.5mm,after skip=12pt]
{\sffamily\bfseries\fontsize{8}{10}\selectfont\textcolor{milcNegro!55}{ESTA GUÍA ES DE}}\\[3mm]
\begin{tabularx}{\linewidth}{X p{3.2cm} p{3.2cm}}
\linead{\linewidth} & \linead{3.0cm} & \linead{3.0cm}\\[-1mm]
{\fontsize{8.5}{10}\selectfont\textcolor{milcNegro!55}{Nombre completo}} &
{\fontsize{8.5}{10}\selectfont\textcolor{milcNegro!55}{Grupo}} &
{\fontsize{8.5}{10}\selectfont\textcolor{milcNegro!55}{Fecha}}\\
\end{tabularx}
\end{tcolorbox}

\vfill

\noindent\textcolor{milcLinea}{\rule{\linewidth}{1pt}}\\[2mm]
{\sffamily\bfseries\fontsize{9.5}{12}\selectfont Autor: Dr. Álvaro Cárdenas Orozco}\\[1mm]
{\sffamily\fontsize{8.5}{11}\selectfont\textcolor{milcNegro!55}{Metodología MILC · apertura ancestral · pensamiento computacional · triángulo Dussel–estoicismo–Floridi}}

\newpage

\section*{Apertura: ¿Qué es un algoritmo? --- la receta paso a paso}

\begin{softbox}{milcGradeDark}{milcGradeSoft}{Saber ancestral}
\textbf{Doña Mercedes la maestra rural de Cartago siempre decía: ``El que se sabe la receta exacta, hornea pan. El que la sabe a medias, hornea decepción.''} En la vereda La Plata, las recetas de panadería se transmitían \emph{exactas, paso a paso, sin saltarse nada}. Si una abuela enseñaba a su nieta a hacer pan de yuca, la receta era: \emph{``1 libra de yuca rallada. 1 huevo. 200 gramos de queso. 1 cucharadita de sal. (1) Pelar y rallar la yuca. (2) Mezclar la yuca con el huevo. (3) Agregar el queso rallado. (4) Amasar con las manos hasta que quede uniforme. (5) Formar bolitas del tamaño de un huevo. (6) Hornear a 180° por 20 minutos. (7) Sacar y dejar enfriar.''}. La nieta podía hornear el mismo pan que su abuela porque la receta era \textbf{exacta, completa y ordenada}. Si alguien decía \emph{``hornee como Dios le ayude''}, la nieta dañaba el pan. \textbf{Una receta clara es un algoritmo.} Hoy en programación llamamos \emph{algoritmo} exactamente a eso: una receta paso a paso para resolver un problema, lo suficientemente clara para que cualquiera (o un computador) pueda ejecutarla.
\end{softbox}

\begin{softbox}{milcTurquesa}{milcGris}{Saber contemporáneo}
Un \textbf{algoritmo} es una \emph{secuencia ordenada de pasos para resolver un problema o realizar una tarea}. La palabra viene del matemático persa \textbf{Al-Juarismi} (siglo IX), que enseñó a Europa los procedimientos matemáticos sistemáticos. Hoy los algoritmos están \textbf{en todas partes}: Google los usa para buscar; TikTok los usa para recomendarte videos; el cajero del banco los usa para entregar dinero; el GPS los usa para encontrar rutas. \textbf{Un buen algoritmo tiene 5 características} (las aprenderás hoy). Si falta una, el algoritmo no funciona bien. \textbf{La habilidad de escribir buenos algoritmos} es el corazón del pensamiento computacional. No requiere computador: puedes escribir un algoritmo en cuaderno, ejecutarlo tú mismo, verificarlo. Hoy haces eso: tu primer algoritmo en lenguaje natural.
\end{softbox}

\begin{softbox}{milcMostaza}{milcGris}{Pregunta-puente}
\textit{Si tu mejor amigo nunca ha hecho una arepa, ¿\textbf{qué tan detallada} debería ser tu receta para que él la haga bien? ¿Cuántos pasos crees que necesita?}
\end{softbox}

\begin{softbox}{milcVerde}{milcGris}{Saber y saber hacer}
Al terminar la clase: (1) podrás \textbf{identificar} qué es un algoritmo y dónde se usa; (2) sabrás \textbf{explicar} las 5 características de un buen algoritmo; (3) podrás \textbf{aplicar} la estructura entrada-proceso-salida; (4) habrás \textbf{creado} tu primer algoritmo en lenguaje natural.
\end{softbox}



\small
\begin{tabularx}{\textwidth}{>{\bfseries}p{3.0cm}Y}
\rowcolor{milcGradePrimary!12}Autor & Dr. Álvaro Cárdenas Orozco\\
DBA & Construye algoritmos y diagramas de flujo para resolver problemas paso a paso (MEN, T\&I 7°)\\
Referentes & Saber ancestral del tejedor · Pensamiento computacional · Scratch\\
Producto final & Un \textbf{algoritmo escrito en el cuaderno} de algo cotidiano que tú sepas hacer (preparar arepas, organizar el morral, lavarte los dientes, hacer una llamada, jugar un juego, etc.). El algoritmo tiene: \textbf{(a) nombre claro}; \textbf{(b) entrada} (qué necesitas para empezar); \textbf{(c) pasos numerados} (mínimo 6); \textbf{(d) salida} (qué resultado se obtiene); \textbf{(e) 5 características} verificadas (claro, finito, determinista, eficaz, ordenado). Más \textbf{una autoevaluación} de cuál característica te costó más cumplir.\\
\end{tabularx}
\normalsize

\puente{Hoy haces 4 cosas: descubres qué es un algoritmo, aprendes las 5 características, escribes tu primer algoritmo, lo verificas.}

\milcestacion{milcGradeDark}{Ruta MILC: así vamos a aprender}
\begin{tabularx}{\textwidth}{>{\centering\arraybackslash}X>{\centering\arraybackslash}X>{\centering\arraybackslash}X>{\centering\arraybackslash}X}
\phasecard{1}{milcVerde}{milcGris}{Escucha\\[-1mm]\small Observo, escucho y pregunto.} &
\phasecard{2}{milcTurquesa}{milcGris}{Entender\\[-1mm]\small Sistematizo y ordeno.} &
\phasecard{3}{milcMagenta}{milcGris}{Hacer\\[-1mm]\small Construyo y aplico.} &
\phasecard{4}{milcVino}{milcGris}{Cierre\\[-1mm]\small Evalúo y reflexiono.}
\end{tabularx}


\puente{Empezamos analizando una receta mala y una buena.}

\milcestacion{milcVerde}{Estación 1 de 3 · Escucha}

\begin{softbox}{milcVerde}{milcGris}{Escena para escuchar}
\textbf{Actividad 1 · ANALIZA --- 2 recetas: ¿cuál sí funciona?} (10 min · individual). Lee estas 2 versiones de \emph{cómo lavar las manos}. En el cuaderno responde: \emph{¿cuál sirve como algoritmo y cuál no?, ¿por qué?}. \textbf{Versión A:} \emph{``Lavarse las manos: agua y jabón. Mojarlas, jabonarlas, frotar, secar.''}. \textbf{Versión B:} \emph{``Lavarse las manos: (1) Abrir el grifo del lavamanos. (2) Mojarse las manos con agua. (3) Poner jabón en la palma de una mano (medio tapón). (4) Frotar las dos manos durante 20 segundos cubriendo dedos, palmas y dorso. (5) Enjuagar con agua hasta que no quede jabón. (6) Cerrar el grifo. (7) Secarse con una toalla limpia.''}. \emph{Ambas describen la misma tarea, pero una sirve como algoritmo y otra no.}
\end{softbox}

{\sffamily\bfseries\fontsize{8}{10}\selectfont\textcolor{milcNegro!55}{MARCA CUANDO LO HAYAS HECHO}}
\milccheckrow{Leí las 2 versiones con calma.}
\milccheckrow{Identifiqué cuál sí sirve como algoritmo.}
\milccheckrow{Pensé qué hace diferente a la versión B.}

\begin{infoband}{milcVerde}


\textbf{Lo que va al cuaderno (Actividad 1 · ANALIZA).} Título: «Actividad 1 --- 2 recetas comparadas». \textbf{Formato:} respuesta corta de 3-4 líneas. \textbf{Extensión:} 4 líneas. \textbf{Sabes que terminaste cuando} identificas que la versión B es algoritmo porque es ordenada, exacta y completa.
\end{infoband}

\puente{Después aprendes las 5 características + la estructura entrada-proceso-salida.}

\milcestacion{milcTurquesa}{Estación 2 de 3 · Entender}

\begin{softbox}{milcTurquesa}{milcGris}{Lo que vas a entender}
La diferencia entre la Versión A y la Versión B es que B cumple las \textbf{5 características de un buen algoritmo}. Ahora las aprendes con detalle.
\end{softbox}

\milcpaso{1}{milcTurquesa}{\textbf{Característica 1 --- CLARO.} Cada paso se entiende sin ambigüedad. \emph{``Mojarlas''} no es claro (¿en qué agua? ¿cuánto?). \emph{``Mojarse las manos con agua del grifo durante 3 segundos''} sí es claro. \textbf{Característica 2 --- FINITO.} El algoritmo tiene un fin. No corre eternamente. Termina en algún paso. \emph{``Lavar manos''} termina al secar. \textbf{Característica 3 --- DETERMINISTA.} Si lo ejecutas con la misma entrada, da el mismo resultado siempre. \emph{Lavar manos con agua y jabón siempre da manos limpias, no a veces sí y a veces no}.}
\milcpaso{2}{milcTurquesa}{\textbf{Característica 4 --- EFICAZ.} Logra el objetivo. No es algoritmo si los pasos no llevan a la meta. \emph{``Bailar mientras mojas las manos''} podría ser claro y finito, pero no eficaz para limpiarlas. \textbf{Característica 5 --- ORDENADO.} Los pasos están en secuencia exacta. \emph{No se puede secar antes de enjuagar; no se puede jabonar antes de mojar}. El orden importa. Si cambias el orden, el algoritmo falla o se vuelve absurdo.}
\milcpaso{3}{milcTurquesa}{\textbf{La estructura entrada-proceso-salida.} Todo algoritmo tiene 3 partes: \textbf{(1) Entrada (input):} lo que necesitas para empezar. En la receta de pan: \emph{yuca, huevo, queso, sal}. \textbf{(2) Proceso:} los pasos ordenados que transforman la entrada. \emph{Mezclar, amasar, formar, hornear}. \textbf{(3) Salida (output):} el resultado final. \emph{Pan de yuca terminado}. Esa estructura aplica a cualquier algoritmo: receta, programa, manual de armado.}
\milcpaso{4}{milcTurquesa}{\textbf{4 ejemplos de algoritmos cotidianos.} (1) \emph{Cepillarse los dientes}: entrada = cepillo + crema + agua; proceso = mojar cepillo, poner crema, cepillar 2 min, enjuagar; salida = boca limpia. (2) \emph{Hacer una llamada}: entrada = celular + número; proceso = desbloquear, ir a Teléfono, marcar número, pulsar verde, hablar, colgar; salida = conversación realizada. (3) \emph{Llegar al colegio}: entrada = uniforme + morral + transporte; proceso = vestirse, desayunar, salir, tomar bus, bajar en la parada, caminar, entrar; salida = llegada a tiempo. (4) \emph{Resolver una ecuación}: entrada = ecuación; proceso = pasar términos, simplificar, despejar; salida = valor de x.}

\begin{drawbox}{milcTurquesa}{5 características + entrada-proceso-salida}
\textbf{5 características:} \\[1mm]
\textbf{1.} CLARO --- cada paso se entiende. \\[1mm]
\textbf{2.} FINITO --- termina en algún momento. \\[1mm]
\textbf{3.} DETERMINISTA --- mismo input, mismo output. \\[1mm]
\textbf{4.} EFICAZ --- logra el objetivo. \\[1mm]
\textbf{5.} ORDENADO --- pasos en secuencia exacta. \\[1mm]
\textbf{Estructura:} ENTRADA → PROCESO → SALIDA.
\end{drawbox}

\begin{softbox}{milcMostaza}{milcGris}{Errores comunes y su corrección}
\textbf{Errores al escribir tu primer algoritmo.} (1) \emph{Pasos vagos:} \emph{``cocinar bien''} no es claro. Usa verbos exactos. (2) \emph{Saltar pasos obvios:} para ti es obvio que primero hay que abrir el grifo, pero para un robot no. Escribe TODO. (3) \emph{Orden incorrecto:} secar antes de enjuagar = manos con jabón. El orden importa. (4) \emph{Sin definir la entrada:} si no dices qué se necesita, quien ejecuta no sabe por dónde empezar. (5) \emph{Sin definir la salida:} si no dices qué resultado esperar, no se sabe si funcionó.
\end{softbox}

\begin{infoband}{milcTurquesa}


\textbf{Lo que va al cuaderno (Actividad 2 · EVALÚA).} Título: «Actividad 2 --- Diagnóstico de un algoritmo defectuoso». \textbf{Formato:} tabla de 5 filas, 2 columnas. Columna 1: característica. Columna 2: ejemplo bueno + ejemplo malo. \textbf{Veredicto:} mira este algoritmo de \emph{``preparar un sandwich''} → ``1. tomar pan 2. tomar jamón 3. ya está''. ¿\textbf{Cuáles de las 5 características le faltan} y cuál es la más grave? Reescríbelo aplicando solo la característica más crítica. \textbf{Extensión:} 5 filas + veredicto. \textbf{Sabes que terminaste cuando} puedes señalar con precisión qué falla y por qué falla en cualquier algoritmo defectuoso.
\end{infoband}

\puente{Con todo claro, escribes tu propio algoritmo paso a paso.}

\milcestacion{milcMagenta}{Estación 3 de 3 · Hacer}

\begin{softbox}{milcMagenta}{milcGris}{Cómo vas a construir tu producto}
\textbf{Actividad 3 · CREA --- Tu primer algoritmo escrito} (28 min · individual). Vas a escribir un algoritmo de algo que tú sepas hacer y que sea cotidiano (no académico). Después lo verificas con un compañero.
\end{softbox}

\milcpaso{1}{milcMagenta}{\textbf{Paso 1 --- Escoge un tema cotidiano.} En el cuaderno escribe el tema. Sugerencias: \emph{``Cómo preparar arepas''}, \emph{``Cómo organizar mi morral''}, \emph{``Cómo enviar un mensaje de WhatsApp''}, \emph{``Cómo jugar [tu juego favorito]''}, \emph{``Cómo hacer una llamada''}, \emph{``Cómo armar un rompecabezas''}. Escoge algo que TÚ sepas hacer bien.}
\milcpaso{2}{milcMagenta}{\textbf{Paso 2 --- Define entrada y salida.} Arriba del algoritmo escribe: \emph{``Algoritmo: [tu tema]''}. Debajo: \emph{``Entrada (lo que necesito): \_\_\_\_\_''} (lista de ingredientes/herramientas/datos). Después: \emph{``Salida (resultado esperado): \_\_\_\_\_''} (qué se obtiene al final).}
\milcpaso{3}{milcMagenta}{\textbf{Paso 3 --- Lista numerada de pasos.} Escribe mínimo \textbf{6 pasos} numerados (1, 2, 3, ...). Cada paso debe ser \emph{una acción concreta} con \emph{verbo en infinitivo}: \emph{``Abrir el grifo''}, \emph{``Mezclar los ingredientes''}, \emph{``Esperar 5 minutos''}. No vagos como \emph{``preparar todo''}. Si te toma 10 pasos, mejor.}
\milcpaso{4}{milcMagenta}{\textbf{Paso 4 --- Verifica las 5 características.} Lee tu algoritmo y marca con palomita ✓ cada una. \emph{¿Es CLARO?} (cada paso se entiende sin ambigüedad). \emph{¿Es FINITO?} (termina, no corre eternamente). \emph{¿Es DETERMINISTA?} (mismo input, mismo output). \emph{¿Es EFICAZ?} (logra el objetivo). \emph{¿Es ORDENADO?} (los pasos están en secuencia correcta). Si alguna falla, corrige antes de seguir.}
\milcpaso{5}{milcMagenta}{\textbf{Paso 5 --- Verifica con un compañero.} Intercambia tu cuaderno con un compañero. Cada uno lee el algoritmo del otro y trata de ejecutarlo mentalmente (si fuera un robot, ¿podría hacerlo?). Anota \emph{dudas o ambigüedades} que encontraste en el algoritmo del compañero. Devuélvanse los cuadernos y corrijan según las dudas señaladas.}
\milcpaso{6}{milcMagenta}{\textbf{Paso 6 --- Autoevaluación.} Al final escribe: \emph{``La característica que más me costó cumplir fue: \_\_\_\_\_ porque \_\_\_\_\_''}. Firma con nombre y fecha. Tu primer algoritmo está listo.}

\begin{drawbox}{milcMagenta}{Antes de cerrar la S2}
\checkbox Mi algoritmo tiene nombre claro arriba. \\[1mm]
\checkbox Entrada y salida están definidas. \\[1mm]
\checkbox Hay mínimo 6 pasos numerados con verbos concretos. \\[1mm]
\checkbox Verifiqué las 5 características y todas se cumplen. \\[1mm]
\checkbox Un compañero lo verificó y corregí lo que él señaló. \\[1mm]
\checkbox Hice autoevaluación firmada.
\end{drawbox}

\begin{softbox}{milcMagenta}{milcGris}{Plantilla de trabajo}
\textbf{Estructura del algoritmo.} \textit{Título:} Algoritmo + tema. \textit{Entrada:} lista de necesidades. \textit{Salida:} resultado esperado. \textit{Pasos:} lista numerada (mínimo 6), cada uno con verbo concreto. \textit{Verificación:} 5 características con palomita. \textit{Autoevaluación:} cuál característica costó más + firma.
\end{softbox}

\begin{infoband}{milcMagenta}


\textbf{Lo que va al cuaderno (Actividad 3 · CREA).} Título: «Actividad 3 --- Mi primer algoritmo». \textbf{Formato:} título + entrada + 6+ pasos + salida + 5 características verificadas + autoevaluación. \textbf{Extensión:} 1 página. \textbf{Sabes que terminaste cuando} otro compañero puede ejecutar tu algoritmo solo leyéndolo (sin que tú expliques).
\end{infoband}

\puente{El producto es ese algoritmo + autoevaluación.}

\milcestacion{milcMostaza}{Producto de la sesión}
\textbf{Mi primer algoritmo en lenguaje natural · verificado y firmado}.
\begin{softbox}{milcMostaza}{milcGris}{Criterios de calidad}
(1) El algoritmo tiene nombre claro y estructura entrada-proceso-salida. (2) Hay mínimo 6 pasos numerados con verbos concretos. (3) Las 5 características están verificadas con palomita. (4) Un compañero verificó el algoritmo y se hicieron correcciones. (5) Hay autoevaluación honesta con firma.
\end{softbox}

\puente{Antes de salir, verifica con un compañero que él podría ejecutar tu algoritmo solo leyéndolo.}

\milcestacion{milcVino}{Evaluación liberadora · cinco dimensiones}

\begin{softbox}{milcVino}{milcGris}{Sentido de la evaluación}
Evaluar no es solo revisar una nota. Es reconocer qué aprendí, qué cambió en mí, qué emociones manejé, cómo me relaciono con la ciudad y la comunidad, y qué le dejaré a quien viene detrás.
\end{softbox}

\textbf{Mi reflexión liberadora en cinco dimensiones.} Completa cada frase iniciada:

\small
\begin{tabularx}{\textwidth}{>{\bfseries\RaggedRight\arraybackslash}p{3.4cm}Y>{\RaggedRight\arraybackslash}p{4.6cm}}
\rowcolor{milcVino}\color{white}Dimensión & \color{white}\bfseries Mi respuesta (frame starter) & \color{white}\bfseries Algo que descubrí\\
1. Personal & Lo que más cambió en mí fue \linead{6.5cm} & \linead{4.2cm}\\
2. Emocional & La emoción más fuerte fue \linead{2.5cm} y la regulé \linead{3cm} & \linead{4.2cm}\\
3. Ciudadana & En democracia esto importa porque \linead{6cm} & \linead{4.2cm}\\
4. Local (Cartago) & En mi barrio sería útil para \linead{6.5cm} & \linead{4.2cm}\\
5. Intergeneracional & A mi abuela/abuelo le diría \linead{6.5cm} & \linead{4.2cm}\\
\end{tabularx}
\normalsize

\begin{infoband}{milcVino}
\textbf{Para la próxima sustentación.} Vuelve a esta tabla en seis meses. Verás que las respuestas cambian — no porque lo aprendido cambie, sino porque tú cambias. La evaluación liberadora es una conversación contigo a través del tiempo.
\end{infoband}


\puente{Cerramos con 3 voces sobre por qué la claridad en los pasos es virtud práctica.}

\milcestacion{milcGradeDark}{Triángulo de pensamiento · cierre filosófico}

\begin{softbox}{milcVino}{milcGris}{Dussel · lente del nosotros}
\textit{El saber que se transmite paso a paso es saber que sobrevive generaciones. El que se transmite con vaguedad muere con quien lo tuvo.}

Doña Mercedes y las abuelas tejedoras del Valle del Cauca \textbf{enseñaban en pasos exactos} porque sabían que solo así el saber sobrevivía. La receta de pan de yuca llegó a tu generación porque cada abuela la escribía clara. La receta del tejido Wayuu llegó a 2026 porque cada generación la transmite con precisión. \emph{Escribir buenos algoritmos es continuar esa tradición}: dejar el saber listo para que otros lo ejecuten, incluso si tú ya no estás ahí para explicar.

\textbf{Si yo enseñara algo que sé hacer a un primo más pequeño, ¿lo enseñaría con la claridad de doña Mercedes, o con vaguedad?}
\end{softbox}
\linead{\linewidth}\\[1mm]
\linead{\linewidth}

\begin{softbox}{milcMostaza}{milcGris}{Estoicismo · Marco Aurelio · cuidado interior}
\textit{Pasos en orden, palabras exactas, fin claro. Esa es la trinidad del trabajo bien hecho.}

\textbf{Las 5 características de un buen algoritmo son disciplina trasladada al pensamiento}: \emph{claro (palabras exactas), finito (con fin claro), ordenado (en secuencia), eficaz (logra meta), determinista (mismo input mismo output)}. Esa disciplina es valiosa en programación pero también en \emph{escribir un ensayo, planear un viaje, resolver una pelea}. La capacidad de descomponer en pasos claros es signo de mente entrenada.

\textbf{¿Qué problema actual de mi vida podría resolver mejor si lo pensara como algoritmo?}
\end{softbox}
\linead{\linewidth}\\[1mm]
\linead{\linewidth}

\begin{softbox}{milcTurquesa}{milcGris}{Floridi · lente de la infoesfera}
\textit{Vivimos rodeados de algoritmos invisibles. El que aprende a escribir uno empieza a entender el mundo digital en que vive.}

Los algoritmos que recomiendan tus videos de TikTok, los que clasifican tus correos como spam, los que ordenan los resultados de Google: \textbf{los escribieron seres humanos como tú}. Eran pensamientos paso a paso que alguien tradujo a lenguaje de programación. \emph{Aprender a escribir algoritmos en lenguaje natural es el primer paso} para entender (y eventualmente criticar) los algoritmos invisibles que rigen tu mundo digital.

\textbf{Si los algoritmos los escriben personas, ¿quién decide qué algoritmo me recomienda mis videos? ¿Y por qué me debería importar?}
\end{softbox}
\linead{\linewidth}\\[1mm]
\linead{\linewidth}

\begin{infoband}{milcNegro}
\textbf{Nota para el docente.} Las tres formulaciones del triángulo son la síntesis del autor de la guía, no citas textuales de los pensadores. Si vas a citarlos en clase, remite a la obra original.
\end{infoband}

\milcestacion{milcVino}{Mi autoevaluación · marca una casilla por fila}

{\small
\setlength{\extrarowheight}{2.2pt}
\begin{tabularx}{\textwidth}{>{\RaggedRight\arraybackslash\bfseries}p{3.0cm}YYY}
\rowcolor{milcVino}\color{white}\bfseries Criterio & \color{white}\bfseries Lo logré & \color{white}\bfseries En proceso & \color{white}\bfseries Necesito apoyo\\[.6mm]
Comprensión del tema & \milcopcion{Explico las ideas con mis palabras.} & \milcopcion{Reconozco las ideas, falta precisión.} & \milcopcion{Necesito repasar lo básico.}\\[.8mm]
\rowcolor{milcGris}Pensamiento computacional & \milcopcion{Usé los pilares al armar mi producto.} & \milcopcion{Usé algunos pilares.} & \milcopcion{Necesito releer la estación 2.}\\[.8mm]
Aplicación y producto & \milcopcion{Mi producto cumple los criterios.} & \milcopcion{Está armado, con ajustes pendientes.} & \milcopcion{Aún está en construcción.}\\[.8mm]
\rowcolor{milcGris}Reflexión 5 dimensiones & \milcopcion{Respondí cada una con honestidad.} & \milcopcion{Respondí algunas con detalle.} & \milcopcion{Mis respuestas son muy generales.}\\[.8mm]
Triángulo de pensamiento & \milcopcion{Conecté las 3 voces con mi trabajo.} & \milcopcion{Conecté al menos una.} & \milcopcion{No las relaciono con mi proceso.}\\[.8mm]
\end{tabularx}}

\vspace{3mm}
\textbf{Hoy entendí que:}\\[1mm]
\linead{\linewidth}\\[1mm]
\linead{\linewidth}

\textbf{Mi compromiso para la próxima sesión es:}\\[1mm]
\linead{\linewidth}\\[1mm]
\linead{\linewidth}

\begin{softbox}{milcMostaza}{milcGris}{Cita de despedida · Marco Aurelio}
\textit{``El obstáculo en el camino se vuelve el camino. Lo que estorba al camino se vuelve camino.''}\\[1mm]
Cada error de hoy, bien observado, se convierte en pieza del camino que pisarás mañana.
\end{softbox}

\small
\begin{tabularx}{\textwidth}{>{\bfseries}p{3.6cm}Y}
\rowcolor{milcGradeDark!12}Nombre completo: & \linead{10cm}\\
Fecha: & \linead{4cm} \quad Curso/grupo: \linead{4cm}\\
Mi firma: & \linead{9cm}\\
\end{tabularx}
\normalsize

\begin{infoband}{milcGradeDark}
\textbf{Para la próxima sesión.} Esta semana, escribe en tu cuaderno un \emph{segundo} algoritmo de otra tarea cotidiana. Practica las 5 características. La próxima clase aprendes las \textbf{4 técnicas del pensamiento computacional}: cómo abordar problemas grandes descomponiéndolos.
\end{infoband}

\end{document}
